\chapter*{Sommario}
\addcontentsline{toc}{chapter}{Sommario}

La crescente integrazione di infrastrutture satellitari all'interno delle moderne reti di comunicazione introduce nuove esigenze in termini di monitoraggio e rilevamento delle intrusioni. 
In tale contesto, l'impiego di sistemi di \textit{Network Intrusion Detection} basati su tecniche di \textit{Machine Learning} risulta particolarmente promettente, ma è condizionato dalla capacità dei modelli di mantenere prestazioni adeguate quando la distribuzione dei dati osservata in fase operativa differisce da quella utilizzata durante l'addestramento.

Il presente lavoro analizza il problema della generalizzazione \textit{cross-domain} di classificatori addestrati su traffico di rete terrestre e successivamente valutati su traffico riconducibile al dominio satellitare. 
A tale scopo è stato sviluppato un framework metodologico che armonizza i dataset UNSW-NB15 e STIN all'interno di uno spazio condiviso delle feature mediante un operatore di mappatura volto a preservare la coerenza semantica e dimensionale delle grandezze considerate.

Il protocollo sperimentale confronta tre differenti paradigmi di addestramento: modelli addestrati esclusivamente sul dominio sorgente, una configurazione di baseline denominata \texttt{INJECTION}, nella quale una quota limitata di osservazioni malevole target viene introdotta nel training set sorgente, e modelli ibridi addestrati utilizzando direttamente la componente malevola del dominio target. 
Le sperimentazioni sono state condotte mediante Decision Tree, Random Forest e Histogram-based Gradient Boosting, adottando una procedura multi-seed e una valutazione sistematica sia sul dominio sorgente sia sui differenti benchmark target.

L'analisi quantitativa è stata affiancata da tecniche diagnostiche e di spiegabilità. 
Le proiezioni PCA e le stime di densità KDE sono state utilizzate per caratterizzare le differenze distributive tra i domini, mentre l'analisi SHAP ha consentito di esaminare la rilevanza delle feature nelle differenti configurazioni. 
La robustezza delle evidenze sperimentali è stata inoltre approfondita attraverso l'analisi della variabilità inter-seed e l'impiego del test $t$ di Welch.

I risultati mostrano una marcata differenza tra i paradigmi considerati.
I modelli addestrati esclusivamente sul dominio sorgente mantengono buone prestazioni sui dati omologhi, ma evidenziano una significativa riduzione della sensibilità nella valutazione \textit{cross-domain}. 
I modelli ibridi raggiungono invece prestazioni molto elevate sui benchmark target, a fronte di una perdita della capacità di riconoscere gli attacchi appartenenti al dominio sorgente. 
La baseline \texttt{INJECTION} presenta il comportamento più equilibrato, preservando sostanzialmente le prestazioni sul dominio sorgente e incrementando al contempo la capacità di rilevamento sul dominio target.

Nel complesso, il lavoro evidenzia come la composizione del dataset di addestramento rappresenti un elemento determinante per la generalizzazione di un NIDS in presenza di \textit{domain shift}. 
I risultati ottenuti mostrano inoltre l'importanza di affiancare alla valutazione in-domain una sperimentazione \textit{cross-domain}, al fine di caratterizzare in modo più realistico la capacità di generalizzazione dei modelli destinati a scenari di rete eterogenei.
